\scanpage{64}
\mspage{63}
\mspage{34}

de caractéristique $p$ avec $(n,p)=1$. Alors $X$ est régulier et chaque faisceau inversible sur un ouvert $j:U\longrightarrow X$ est restriction d'un faisceau sur $X$, d'où $R^1j_*\mathbb G_m=0$. Il s'ensuite de la suite exacte de Kummer (IX.3.2) qu'on a un isomorphisme canonique
\diagram{\[
(*)\qquad\begin{tikzcd}
(j_*\mathbb G_m)/n\arrow[r,"\sim"]&R^1j_*\mu_n
\end{tikzcd}
\]}
(notation de (3.1)). De plus, on a la suite exacte canonique
\diagram{\[
\begin{tikzcd}
0\arrow[r]&(\mathbb G_m)_X\arrow[r]&j_*(\mathbb G_m)_U\arrow[r]&\mathbb Z_Y\arrow[r]&0
\end{tikzcd}
\]}
où la section unité de \deleted{\unsure{$(\mathbb G_m)$}} $\mathbb Z_Y$ est l'image de la section $t$ de $j_*(\mathbb G_m)_U$ ($t=0$ une équation locale pour $Y$ arbitraire). On en déduit un isomorphisme canonique
\diagram{\[
\begin{tikzcd}
(j_*\mathbb G_m)/n\arrow[r,"\sim"]&\deleted{\unsure{0|\mathbb G}}\ (\mathbb Z/n)_Y
\end{tikzcd}
\]}
\deleted{\unsure{qui}} ce qui, ajouté à (*), donne le résultat.

Supposons maintenant que l'existence et unicité sont déjà démontrés si la codimension est $<c$, et que $c>1$. En appliquant (3.8) et l'hypothèse de recurrence, on constate immédiatement que dans la situation de (3.10)(iii) les trois flèches de \deleted{(3.10.1)} (3.10.1) déjà définis, i.e., $\varphi(Z,Y)$, $\xi$, et $\eta$ sont de nature locale pour la topologie étale et commutent aux changement de base
